% =============================================================
%  Teaching a Network to Bend: Hidden Layers and Activation Functions
%  Series A: Machine Learning Basics for Chemists, Post 2
%  Compiles with pdfLaTeX on Overleaf (no extra files needed).
% =============================================================
\documentclass[11pt]{article}

\usepackage[margin=1in]{geometry}
\usepackage[T1]{fontenc}
\usepackage{lmodern}
\usepackage{microtype}
\usepackage{amsmath,amssymb,bm}
\usepackage{booktabs}
\usepackage{enumitem}
\usepackage{xcolor}
\usepackage{listings}
\usepackage[most]{tcolorbox}
\usepackage[colorlinks=true,linkcolor=blue!50!black,urlcolor=blue!50!black]{hyperref}

\setlength{\parskip}{0.55em}
\setlength{\parindent}{0pt}

% ---------- macros ----------
\newcommand{\T}{\mathsf{T}}
\newcommand{\yhat}{\hat{y}}
\newcommand{\grad}{\nabla}
\DeclareMathOperator{\sign}{sign}
\DeclareMathOperator{\relu}{ReLU}

% ---------- Colab link: edit this ONE line after you upload the notebook to GitHub ----------
\newcommand{\colaburl}{https://colab.research.google.com/github/YOUR-USER/YOUR-REPO/blob/main/notebooks/A2_neural_networks_colab.ipynb}
\newcommand{\colabbutton}{%
  \href{\colaburl}{\tcbox[on line,colback=orange!15,colframe=orange!70!black,boxrule=0.6pt,arc=3pt,
    left=4pt,right=4pt,top=1pt,bottom=1pt]{\small\textbf{$\blacktriangleright$ Open in Google Colab}}}}
\newcommand{\tryit}[1]{\par\vspace{0.3em}\noindent
  \fcolorbox{orange!70!black}{orange!8}{\parbox{\dimexpr\linewidth-2\fboxsep-2\fboxrule\relax}{%
  \small $\blacktriangleright$ \textbf{Run it.} #1}}\par\vspace{0.3em}}

% ---------- code listings ----------
\definecolor{codebg}{RGB}{246,246,248}
\definecolor{codecomment}{RGB}{90,110,90}
\definecolor{codekw}{RGB}{30,60,160}
\lstset{
  language=Python,
  basicstyle=\ttfamily\small,
  keywordstyle=\color{codekw}\bfseries,
  commentstyle=\color{codecomment}\itshape,
  backgroundcolor=\color{codebg},
  frame=single,
  rulecolor=\color{black!15},
  breaklines=true,
  columns=fullflexible,
  showstringspaces=false,
  xleftmargin=0.5em,
  xrightmargin=0.5em,
  aboveskip=0.8em,
  belowskip=0.8em
}

% ---------- boxes ----------
\tcbuselibrary{skins,breakable}
\newtcolorbox{keybox}[1][]{
  enhanced, breakable, colback=blue!4, colframe=blue!50!black,
  fonttitle=\bfseries, title={#1}, boxrule=0.6pt, arc=2pt
}
\newtcolorbox{mathbox}[1][]{
  enhanced, breakable, colback=black!3, colframe=black!45,
  fonttitle=\bfseries, title={#1}, boxrule=0.6pt, arc=2pt
}

% ---------- title ----------
\title{\textbf{Teaching a Network to Bend:}\\[0.2em]
       Hidden Layers and Activation Functions}
\author{\textit{Series A: Machine Learning Basics for Chemists --- Post 2 of 7}}
\date{}

\begin{document}
\maketitle
\thispagestyle{empty}

\begin{center}
\small
About 20 minutes of reading, plus the notebook.\\
\textbf{Prerequisites:} Post A1 (gradient descent) and basic Python. The math supplement covers the rest.\\
\textbf{Companion notebook:} \texttt{A2\_neural\_networks\_colab.ipynb}\\[0.4em]
\colabbutton\quad {\footnotesize (no installation; see Section~\ref{sec:runit})}
\end{center}

\vspace{0.5em}
\hrule
\vspace{0.5em}

% =============================================================
\section{The question}
% =============================================================
In Post A1 we fitted $\sin(x)$ with a cubic polynomial. That worked because we chose the features $[1,z,z^2,z^3]$ ourselves, and only the four coefficients were learned. The best cubic on our 200-point grid still leaves a mean squared error of about $4.6\times10^{-3}$, and it falls apart badly outside the fitted range.

This post asks a bigger question: \textbf{can the model choose its own features?} A neural network does exactly that. It learns both the building blocks and how to combine them. We will build a small one from scratch, write the training algorithm ourselves, and discover that one design choice, the \emph{activation function}, decides whether the whole thing is powerful or useless.

% =============================================================
\section{Step 1: A neuron and its activation function}
% =============================================================
A single \textbf{neuron} takes the input $x$, computes a weighted sum, and passes it through a fixed function $\phi$ called the \textbf{activation}:
\begin{equation}
  z = w\,x + b, \qquad a = \phi(z).
\end{equation}
The weight $w$ and bias $b$ are learned. The activation $\phi$ is chosen by us. Three common choices, plus the ``linear'' one that does nothing:

\begin{center}
\small
\begin{tabular}{@{}lccc@{}}
\toprule
Name & $\phi(z)$ & Output range & Largest $\phi'(z)$\\
\midrule
Linear & $z$ & $(-\infty,\infty)$ & $1$\\
Logistic sigmoid & $1/(1+e^{-z})$ & $(0,1)$ & $1/4$\\
Hyperbolic tangent & $\tanh z$ & $(-1,1)$ & $1$\\
ReLU & $\max(0,z)$ & $[0,\infty)$ & $1$ (for $z>0$)\\
\bottomrule
\end{tabular}
\end{center}

The sigmoid and tanh are smooth S-shaped curves that flatten out for large $|z|$. ReLU (``rectified linear unit'') is zero for negative $z$ and a straight line for positive $z$, so it has one bend, at $z=0$.

\tryit{Notebook \S1. Plot all three activations and their derivatives, and confirm the maximum derivative of each. The sigmoid never exceeds $0.25$.}

% =============================================================
\section{Step 2: A hidden layer}
% =============================================================
Now use $H$ neurons side by side, all looking at the same input, and add up their outputs with learned weights $v_h$:
\begin{equation}
  \yhat(x)=\sum_{h=1}^{H} v_h\,\phi\!\left(w_h x+b_h\right)+c .
  \label{eq:net}
\end{equation}
The neurons in the middle are called the \textbf{hidden layer}, because we never see their outputs directly. Equation~\eqref{eq:net} has three numbers per neuron plus one overall offset $c$, so $3H+1$ parameters in total. With $H=10$ that is $31$, compared to $4$ for the cubic.

Compare this with Post A1. The cubic used fixed building blocks $1,z,z^2,z^3$ and learned how much of each to use. Here the building blocks $\phi(w_hx+b_h)$ are themselves adjustable: $w_h$ controls how steep each one is, and $b_h$ controls where it sits. The network learns a tailor-made set of features.

We use the same data as before, now $N=200$ points of $\sin x$ on $[-\pi,\pi]$, and the same loss without the ridge penalty:
\begin{equation}
  L=\frac1N\sum_{i=1}^N\left(\yhat(x_i)-y_i\right)^2 .
\end{equation}

\tryit{Notebook \S0. The cell defines the whole network (initialization, forward pass, gradient, optimizer) in about 40 lines of NumPy.}

% =============================================================
\section{Step 3: Why stacking linear layers is useless}
% =============================================================
Before trusting any nonlinear activation, it is worth seeing why we need one. Suppose $\phi(z)=z$. Then Eq.~\eqref{eq:net} becomes
\[
  \yhat=\sum_h v_h(w_hx+b_h)+c=\Big(\sum_h v_hw_h\Big)x+\Big(\sum_h v_hb_h+c\Big),
\]
a straight line, no matter how large $H$ is. Adding more linear layers does not help either (\hyperref[sec:S1]{Supplement S1}). All $31$ parameters collapse into one slope and one intercept.

The notebook confirms this. Training a 10-unit linear network on $\sin x$ gives:

\begin{center}
\begin{tabular}{@{}lccc@{}}
\toprule
 & Slope & Intercept & MSE\\
\midrule
Best straight line & $0.2994$ & $0$ & $0.1996$\\
10-unit linear network & $0.2994$ & $0$ & $0.1996$\\
\bottomrule
\end{tabular}
\end{center}

They are identical. The depth bought nothing. A hidden layer is only worth having if the activation is \textbf{nonlinear}.

\tryit{Notebook \S2. Train the linear network and compute its collapsed slope as $\sum_h w_hv_h$.}

% =============================================================
\section{Step 4: Backpropagation, the gradient of the network}
% =============================================================
To train the network we need the gradient of $L$ with respect to all $3H+1$ parameters, and this is where the neural network earns its famous algorithm. It is the chain rule, organized from the output backwards.

Write $z_{ih}=w_hx_i+b_h$, $a_{ih}=\phi(z_{ih})$, and let $r_i=\yhat_i-y_i$ be the residual. Then (derived in \hyperref[sec:S3]{Supplement S3}):

\begin{align}
  \delta^{(2)}_i &= \tfrac{2}{N}\,r_i
  &&\text{error at the output}\notag\\
  \partial L/\partial v_h &= \sum_i a_{ih}\,\delta^{(2)}_i,
  \qquad
  \partial L/\partial c = \sum_i \delta^{(2)}_i
  &&\text{output layer}\notag\\
  \delta^{(1)}_{ih} &= \delta^{(2)}_i\,v_h\,\phi'(z_{ih})
  &&\text{error pushed back through }\phi\notag\\
  \partial L/\partial w_h &= \sum_i x_i\,\delta^{(1)}_{ih},
  \qquad
  \partial L/\partial b_h = \sum_i \delta^{(1)}_{ih}
  &&\text{hidden layer}\notag
\end{align}

Notice where the activation appears: the error is multiplied by $\phi'(z_{ih})$ on the way back. If $\phi'$ is tiny, little signal reaches the hidden weights. We will return to this in the next post, on vanishing gradients.

As in Post A1, we test the formulas before trusting them. Comparing against central finite differences, the largest disagreement over all parameters is about $3\times10^{-10}$ for tanh, $3\times10^{-10}$ for sigmoid, and $5\times10^{-10}$ for ReLU.

\tryit{Notebook \S3. Run the finite-difference check for all three activations. You wrote the same check for the cubic in Post A1.}

% =============================================================
\section{Step 5: Train a tanh network}
% =============================================================
We train a 10-unit tanh network with \textbf{Adam}, the optimizer we met at the end of Post A1, using learning rate $0.01$ for $4000$ updates. Starting from random weights (seed 0), the loss evolves like this:

\begin{center}
\begin{tabular}{@{}rc@{}}
\toprule
Updates & Loss (MSE)\\
\midrule
0 & $1.56$\\
10 & $0.50$\\
100 & $5.9\times10^{-2}$\\
500 & $1.7\times10^{-4}$\\
1000 & $1.5\times10^{-5}$\\
2000 & $6.6\times10^{-6}$\\
4000 & $4.0\times10^{-6}$\\
\bottomrule
\end{tabular}
\end{center}

The final error is about $1150$ times smaller than the best cubic, and the largest error anywhere on the grid is $0.0085$. Plain gradient descent with a fixed step size reaches only about $2\times10^{-4}$ in the same number of updates, which is why adaptive optimizers are standard for networks.

It is instructive to look inside. The prediction is a sum of ten scaled S-curves $v_h\tanh(w_hx+b_h)$ plus the offset $c$. The notebook plots each of the ten terms and their sum: no single piece looks like a sine wave, but together they reproduce it.

\tryit{Notebook \S4--5. Train the network, print the loss table above, and plot the ten building blocks.}

% =============================================================
\section{Step 6: What each activation learned}
% =============================================================
Now we repeat the experiment for different activations. Each row averages five random starts (seeds 0 to 4), with $H=10$ and $4000$ Adam updates. The last column evaluates the trained network on $[-2\pi,2\pi]$, twice as wide as the training interval.

\begin{center}
\small
\begin{tabular}{@{}lccc@{}}
\toprule
Activation & Train MSE & MSE on fresh points in $[-\pi,\pi]$ & MSE on $[-2\pi,2\pi]$\\
\midrule
Linear & $2.0\times10^{-1}$ & $2.0\times10^{-1}$ & $2.28$\\
Sigmoid & $1.4\times10^{-5}$ & $1.4\times10^{-5}$ & $0.114$\\
Tanh & $6.4\times10^{-6}$ & $6.1\times10^{-6}$ & $0.088$\\
ReLU & $5.4\times10^{-4}$ & $5.3\times10^{-4}$ & $0.493$\\
\midrule
Best cubic (Post A1) & $4.6\times10^{-3}$ & --- & $32.0$\\
Predict $0$ everywhere & --- & --- & $0.500$\\
\bottomrule
\end{tabular}
\end{center}

Three observations:
\begin{enumerate}[nosep]
  \item \textbf{Linear fails, as predicted.} Its error equals that of the best straight line.
  \item \textbf{The smooth activations win inside the range.} Tanh and sigmoid fit a smooth target to about $10^{-5}$. Fresh points do about as well as the training points, so the network is not just memorizing.
  \item \textbf{ReLU is about 80 times less accurate than tanh here.} A ReLU network is a connected set of \emph{straight segments}, and a smooth curve needs many of them. Each unit has one bend at $x_h=-b_h/w_h$, and units whose bend lies outside the data range help little. For seed 0, only $6$ of the $10$ bends fall inside $[-\pi,\pi]$, and the largest error is $0.073$ (\hyperref[sec:S5]{Supplement S5}).
\end{enumerate}

\tryit{Notebook \S6 and \S8. Reproduce the table, then plot the ReLU network with its bends marked. Try to predict which seeds will do best by counting bends inside the range.}

% =============================================================
\section{Step 7: How many hidden units?}
% =============================================================
Mean training MSE over five seeds, for a range of hidden-layer sizes, with the same $4000$ updates for every size:

\begin{center}
\begin{tabular}{@{}rcc@{}}
\toprule
$H$ & Tanh & ReLU\\
\midrule
1  & $6.5\times10^{-2}$ & $3.0\times10^{-1}$\\
2  & $3.3\times10^{-2}$ & $1.5\times10^{-1}$\\
3  & $5.4\times10^{-5}$ & $1.2\times10^{-1}$\\
5  & $1.6\times10^{-5}$ & $3.9\times10^{-3}$\\
10 & $6.4\times10^{-6}$ & $5.4\times10^{-4}$\\
20 & $1.5\times10^{-5}$ & $1.1\times10^{-4}$\\
50 & $2.5\times10^{-6}$ & $7.5\times10^{-4}$\\
\bottomrule
\end{tabular}
\end{center}

The tanh network drops by three orders of magnitude between two and three units: one smooth wave of sine needs only a few S-curves. ReLU needs more units for the same job.

The table is not monotone, though. With $H=20$ the tanh average is \emph{worse} than with $H=10$. The cause is the optimizer, not the model: one of the five runs (seed 2) is slow to converge and dominates the mean. Rerunning the five $H=20$ networks for $20000$ updates instead of $4000$ lowers the mean to $1.8\times10^{-6}$, better than $H=10$. This is the same distinction as in Post A1 between ``the optimizer has not converged'' and ``the model cannot fit.''

\tryit{Notebook \S7 (about 30 seconds). Then change \texttt{steps=4000} to \texttt{steps=20000} for $H=20$ and compare.}

% =============================================================
\section{Step 8: Outside the training range}
% =============================================================
The last column of the Step~6 table is a warning. Inside $[-\pi,\pi]$ the tanh network is excellent. On $[-2\pi,2\pi]$ it does only slightly better than always predicting $0$ would (MSE $0.5$), and the cubic does catastrophically worse ($32$). The network has learned the \emph{shape of the curve on the data}, not the fact that sine repeats.

The behavior outside the data is predictable. A tanh or sigmoid unit flattens to a constant, so the whole network settles to a constant on each side. A ReLU network continues as a straight line (\hyperref[sec:S8]{Supplement S8}). Neither continues the oscillation.

\begin{keybox}[A rule worth remembering]
A trained network is a reliable \emph{interpolator} and an unreliable \emph{extrapolator}. For molecules, this means predictions for structures unlike anything in the training set deserve extra suspicion.
\end{keybox}

\tryit{Notebook \S9. Plot the sine, the cubic, the tanh network and the ReLU network over $[-2\pi,2\pi]$ with the training range shaded.}

% =============================================================
\section{Run it yourself in Colab}
\label{sec:runit}
% =============================================================
Everything above can be reproduced in your browser with no installation. The notebook takes about a minute to run in full.

\begin{keybox}[Three ways to run it]
\begin{enumerate}[nosep]
  \item \textbf{One click.} Press \colabbutton\ (the link points to the notebook in the series repository), then choose \emph{Runtime $\to$ Run all}.
  \item \textbf{Upload the file.} In Colab, choose \emph{File $\to$ Upload notebook} and select \texttt{A2\_neural\_networks\_colab.ipynb}.
  \item \textbf{Paste one cell.} Copy the script in Section~\ref{sec:onecell} into a new Colab cell. It needs only NumPy.
\end{enumerate}
\end{keybox}

\subsection{Where each step lives in the notebook}
\begin{center}
\small
\begin{tabular}{@{}llp{6.8cm}@{}}
\toprule
Post & Notebook & What you should see\\
\midrule
Step 1 & \S1 & Maximum derivatives $0.250$, $1.000$, $1.000$\\
Step 2 & \S0 & Network code; $3H+1=31$ parameters for $H=10$\\
Step 3 & \S2 & Linear net slope $0.2994$, MSE $0.1996$\\
Step 4 & \S3 & Backprop vs.\ finite difference to about $10^{-10}$\\
Step 5 & \S4--5 & Loss table; the cubic is about $1150\times$ worse\\
Step 6 & \S6, \S8 & Activation comparison table; ReLU bends\\
Step 7 & \S7 & The hidden-units table\\
Step 8 & \S9 & Plot showing all models failing outside the range\\
Cross-check & \S10 & PyTorch autograd agrees with our backprop\\
Exercises & \S11 & Starter code for the exercises below\\
Self-check & \S12 & ``All checks passed''\\
\bottomrule
\end{tabular}
\end{center}

\subsection{The core experiment in one cell}
\label{sec:onecell}
This script uses only NumPy. It builds the tanh network, writes out backpropagation, trains with Adam, and compares with the cubic.

\begin{lstlisting}
import numpy as np

N, H = 200, 10
x = np.linspace(-np.pi, np.pi, N); y = np.sin(x); X = x[:, None]
rng = np.random.default_rng(0)
p = {"W1": rng.standard_normal((1, H)), "b1": rng.standard_normal(H),
     "W2": rng.standard_normal((H, 1)) / np.sqrt(H), "b2": np.zeros(1)}

def loss_grad(p):
    A = np.tanh(X @ p["W1"] + p["b1"])               # hidden layer
    r = A @ p["W2"] + p["b2"] - y[:, None]           # residuals
    d2 = 2 * r / N                                   # dL/dyhat
    d1 = (d2 @ p["W2"].T) * (1 - A**2)               # back through tanh
    return np.mean(r**2), {"W2": A.T @ d2, "b2": d2.sum(0), "W1": X.T @ d1, "b1": d1.sum(0)}

m = {k: 0 * v for k, v in p.items()}; s = {k: 0 * v for k, v in p.items()}
for t in range(1, 4001):                             # Adam, lr = 0.01
    L, g = loss_grad(p)
    if t - 1 in (0, 100, 1000): print(f"after {t-1:>4} updates: loss = {L:.3e}")
    for k in p:
        m[k] = 0.9 * m[k] + 0.1 * g[k]; s[k] = 0.999 * s[k] + 0.001 * g[k] ** 2
        p[k] -= 0.01 * (m[k] / (1 - 0.9**t)) / (np.sqrt(s[k] / (1 - 0.999**t)) + 1e-8)
print(f"after 4000 updates: loss = {loss_grad(p)[0]:.3e}")

c = np.polynomial.polynomial.polyfit(x / np.pi, y, 3)          # the cubic from Post 1
print(f"best cubic MSE     : {np.mean((np.polynomial.polynomial.polyval(x/np.pi, c) - y)**2):.3e}")
xf = np.linspace(-2 * np.pi, 2 * np.pi, 2000)                  # outside the training range
net = np.tanh(xf[:, None] @ p["W1"] + p["b1"]) @ p["W2"] + p["b2"]
print(f"MSE on [-2pi, 2pi] : {np.mean((net[:, 0] - np.sin(xf))**2):.3f}   (predicting 0 gives 0.500)")
\end{lstlisting}

Expected output:
\begin{lstlisting}[language={},basicstyle=\ttfamily\small]
after    0 updates: loss = 1.561e+00
after  100 updates: loss = 5.855e-02
after 1000 updates: loss = 1.463e-05
after 4000 updates: loss = 3.958e-06
best cubic MSE     : 4.574e-03
MSE on [-2pi, 2pi] : 0.047   (predicting 0 gives 0.500)
\end{lstlisting}

The last number is for this one network (seed 0). The $0.088$ in the Step~6 table is the average over five seeds.

\subsection{Experiments worth five minutes}
\begin{itemize}[nosep]
  \item \textbf{Swap the activation.} In the one-cell script, replace \texttt{np.tanh} by a sigmoid (and \texttt{1 - A**2} by \texttt{A * (1 - A)}). Does the loss table change much?
  \item \textbf{Change the target.} Replace \texttt{np.sin(x)} by \texttt{np.abs(x)}. Which activation now wins, and why might that be?
  \item \textbf{Shrink the training range.} Train on $[-\pi/2,\pi/2]$ only and evaluate on $[-\pi,\pi]$. How fast does the error grow outside?
  \item \textbf{Break the symmetry on purpose.} Give every hidden unit identical starting weights (Exercise~5) and look at what the ten units do.
\end{itemize}

% =============================================================
\section{What to take away}
% =============================================================
\begin{keybox}[Summary]
\begin{itemize}[nosep]
  \item A hidden layer lets the model \textbf{learn its own features} instead of using ones we choose by hand.
  \item With a \textbf{linear} activation, any number of layers collapses into a straight line. Nonlinearity is not optional.
  \item \textbf{Backpropagation} is the chain rule applied from the output backwards. Each layer multiplies the error by $\phi'$.
  \item Verify gradients with finite differences, as in Post A1.
  \item Smooth activations (tanh, sigmoid) fit smooth targets efficiently. ReLU builds curves from straight pieces and needs more units.
  \item A poor fit can come from an unfinished optimizer and not from the model. Check before changing the architecture.
  \item Networks \textbf{interpolate well and extrapolate badly}.
\end{itemize}
\end{keybox}

\textbf{Next in the series (A3):} why deep networks stop learning. The factor $\phi'$ in the backpropagation formulas, multiplied layer after layer, is the culprit.

\clearpage
% =============================================================
\section*{Supplementary Reading: The Math}
\addcontentsline{toc}{section}{Supplementary Reading: The Math}
% =============================================================
\textit{Prerequisites: matrix multiplication and the chain rule (Toolbox T1, T2). Post A1 for gradient descent and Adam.}

\phantomsection\label{sec:S1}
\subsection*{S1. Linear layers collapse}
Take a network with $K$ layers, all with identity activation, weights $W_1,\dots,W_K$ and biases $\mathbf b_1,\dots,\mathbf b_K$. Applying them in sequence to an input $\mathbf x$:
\[
  \mathbf y=W_K(\cdots(W_2(W_1\mathbf x+\mathbf b_1)+\mathbf b_2)\cdots)+\mathbf b_K
  =\underbrace{(W_K\cdots W_1)}_{W_{\rm eff}}\,\mathbf x+\mathbf b_{\rm eff}.
\]
For two layers, $\mathbf b_{\rm eff}=W_2\mathbf b_1+\mathbf b_2$. Matrix products are matrices, so the whole network is one affine map. Moreover $\operatorname{rank}W_{\rm eff}\le\min_k(\text{width of layer }k)$: extra width cannot add expressiveness either.

For our scalar network with weights $W_1\in\mathbb R^{1\times H}$ and $W_2\in\mathbb R^{H\times1}$:
\[
  \text{slope}=W_1W_2=\sum_h w_hv_h,\qquad
  \text{intercept}=\mathbf b_1^{\T}W_2+c=\sum_h b_hv_h+c .
\]
These are the two numbers computed in the notebook's Section 2.

\subsection*{S2. Forward pass in matrix form}
Stack the $N$ inputs as a column $X\in\mathbb R^{N\times1}$. Then
\[
  Z=XW_1+\mathbf 1\mathbf b_1^{\T}\in\mathbb R^{N\times H},\qquad
  A=\phi(Z)\ \text{(elementwise)},\qquad
  \hat{\mathbf y}=AW_2+c\,\mathbf 1,
\]
and $L=\frac1N\lVert\hat{\mathbf y}-\mathbf y\rVert^2$. For $D$ inputs, $K$ outputs and one hidden layer of width $H$, the parameter count is $(D+1)H+(H+1)K$. Our case $D=K=1$ gives $2H+H+1=3H+1$.

\phantomsection\label{sec:S3}
\subsection*{S3. Deriving backpropagation}
Let $r_i=\hat y_i-y_i$, so $L=\frac1N\sum_ir_i^2$ and
\[
  \delta^{(2)}_i\equiv\frac{\partial L}{\partial\hat y_i}=\frac2N\,r_i .
\]
\textbf{Output layer.} Since $\hat y_i=\sum_ha_{ih}v_h+c$, we have $\partial\hat y_i/\partial v_h=a_{ih}$ and $\partial\hat y_i/\partial c=1$, so
\[
  \frac{\partial L}{\partial v_h}=\sum_i\delta^{(2)}_i\,a_{ih},\qquad
  \frac{\partial L}{\partial c}=\sum_i\delta^{(2)}_i .
\]
In matrix form: $\partial L/\partial W_2=A^{\T}\bm\delta^{(2)}$ and $\partial L/\partial c=\mathbf 1^{\T}\bm\delta^{(2)}$.

\textbf{Hidden layer.} The weight $w_h$ affects $\hat y_i$ only through the single path $w_h\to z_{ih}\to a_{ih}\to\hat y_i$:
\[
  \frac{\partial\hat y_i}{\partial w_h}
  =v_h\,\phi'(z_{ih})\,x_i,
  \qquad
  \frac{\partial\hat y_i}{\partial b_h}
  =v_h\,\phi'(z_{ih}).
\]
Define $\delta^{(1)}_{ih}=\delta^{(2)}_iv_h\phi'(z_{ih})$. Then
\[
  \frac{\partial L}{\partial w_h}=\sum_ix_i\,\delta^{(1)}_{ih},\qquad
  \frac{\partial L}{\partial b_h}=\sum_i\delta^{(1)}_{ih}.
\]
In matrix form, with $\odot$ for the elementwise product:
\begin{mathbox}
\[
  \Delta_1=\big(\bm\delta^{(2)}W_2^{\T}\big)\odot\phi'(Z),\qquad
  \frac{\partial L}{\partial W_1}=X^{\T}\Delta_1,\qquad
  \frac{\partial L}{\partial\mathbf b_1}=\Delta_1^{\T}\mathbf 1 .
\]
\end{mathbox}
This is exactly what the code computes. With more layers the pattern repeats: each layer's $\Delta$ is the next layer's $\Delta$ times a weight matrix, times $\phi'$. The product of many $\phi'$ factors is the subject of Post A3.

\subsection*{S4. Derivatives of the activations}
\paragraph{Sigmoid.} With $\sigma(z)=1/(1+e^{-z})$,
\[
  \sigma'(z)=\frac{e^{-z}}{(1+e^{-z})^2}=\sigma(z)\big(1-\sigma(z)\big)
  =\tfrac14-\big(\sigma(z)-\tfrac12\big)^2\le\tfrac14,
\]
with equality only at $z=0$. In code, $\phi'$ is written as \texttt{a * (1 - a)} using the stored activation $a=\sigma(z)$.

\paragraph{Tanh.} $\tanh z=2\sigma(2z)-1$, so
\[
  \tanh'(z)=4\sigma'(2z)=1-\tanh^2z\in(0,1].
\]
In code: \texttt{1 - a**2}.

\paragraph{ReLU.} $\relu'(z)=1$ for $z>0$ and $0$ for $z<0$. The derivative at $z=0$ does not exist; implementations pick $0$ by convention. A ReLU unit with $z<0$ for \emph{every} training point has zero gradient and never updates, which is a ``dead'' unit.

\phantomsection\label{sec:S5}
\subsection*{S5. ReLU networks are piecewise linear}
In $\hat y(x)=\sum_hv_h\max(0,w_hx+b_h)+c$ each term is linear on either side of its \textbf{kink} $x_h=-b_h/w_h$ (for $w_h\ne0$). The sum is therefore continuous and piecewise linear, with at most $H$ kinks in any interval, hence at most $H+1$ pieces there. A unit whose kink lies outside $[-\pi,\pi]$ is a straight line (or zero) across the whole training range and adds no bend where we need one. In the notebook, seed 0 has $6$ kinks inside the range, seeds 1 to 4 have $9,8,9,5$.

\paragraph{How good can straight pieces be?} Let $f$ satisfy $|f''|\le M$ and let $I_hf$ be its linear interpolant on a mesh of spacing $h$. On one cell $[a,b]$ with $b-a=h$, the standard interpolation error formula gives, for some $\xi\in(a,b)$,
\[
  f(x)-I_hf(x)=\tfrac12f''(\xi)(x-a)(x-b).
\]
Since $\max_{[a,b]}|(x-a)(x-b)|=h^2/4$,
\[
  \max|f-I_hf|\le\frac{M\,h^2}{8}.
\]
For $\sin$ we have $M=1$. With $n$ equal pieces on $[-\pi,\pi]$, $h=2\pi/n$ and the bound is $\pi^2/(2n^2)$. For $n=7$ pieces (the $6$ kinks of seed 0) this is $0.101$. The trained network's largest error is $0.073$, below that bound, which is possible because the training placed the kinks unevenly instead of on an equal mesh.

\subsection*{S6. Universal approximation and how many units are needed}
\paragraph{The theorem.} Cybenko (1989) and Hornik (1991) proved that one-hidden-layer networks with sigmoidal activations can approximate any continuous function on a compact set to any accuracy, if $H$ is allowed to grow. Leshno, Lin, Pinkus and Schocken (1993) showed the same for any continuous activation that is not a polynomial. The polynomial exception is essential: a network with $\phi(z)=z^m$ is itself a polynomial of degree at most $m$.

\paragraph{What the theorem does not say.} It gives no bound on $H$ and no promise that gradient descent will find the weights. The experiments in Step~7 show that both matter.

\paragraph{A concrete count for ReLU.} A continuous piecewise-linear function with $n$ pieces on $[-\pi,\pi]$ can be written as a sum of $n$ ReLU units: one with its kink at the left end to supply the starting slope, and one for each of the $n-1$ interior kinks. Combining with S5, achieving error $\varepsilon$ for $\sin x$ takes at most
\[
  \frac{\pi^2}{2n^2}\le\varepsilon\iff n\ge\frac{\pi}{\sqrt{2\varepsilon}}
\]
units. For $\varepsilon=10^{-3}$ that is $n\ge71$. The trained tanh network with only $10$ units reaches a maximum error of $0.0085$. Smooth activations are far more economical for smooth targets.

\subsection*{S7. Why random initialization is necessary}
Suppose all hidden units start identical: $(w_h,b_h,v_h)=(w,b,v)$ for every $h$. Then $a_{ih}=\phi(wx_i+b)$ is the same for all $h$, so by the formulas in S3,
\[
  \frac{\partial L}{\partial v_h}=\sum_i\delta^{(2)}_ia_{ih},\qquad
  \frac{\partial L}{\partial w_h}=\sum_ix_i\delta^{(2)}_iv\,\phi'(wx_i+b)
\]
are the same for every $h$. Every update changes all units by the same amount, so by induction they stay identical forever. The network is then equivalent to a single unit with output weight $Hv$. In the notebook this gives a final MSE of $0.0655$, close to the $0.0654$ of the $H=1$ network in Step~7.

Random starting weights break this symmetry. We also scale the output weights as $v_h\sim\mathcal N(0,1/H)$: if the $a_h$ are roughly independent with $\mathbb E[a_h^2]=O(1)$, the variance of $\sum_hv_ha_h$ is $\sum_h\frac1H\mathbb E[a_h^2]=O(1)$ independent of $H$, so the starting predictions do not blow up when $H$ grows.

\phantomsection\label{sec:S8}
\subsection*{S8. Behavior outside the training range}
Let $|x|\to\infty$ with $w_h\ne0$.
\begin{itemize}[nosep]
  \item \textbf{Tanh.} $\tanh(w_hx+b_h)\to\sign(w_h)\,\sign(x)$, so
  $\hat y\to c\pm\sum_hv_h\sign(w_h)$, a constant on each side.
  \item \textbf{Sigmoid.} The same, with limits $0$ or $1$ for each unit.
  \item \textbf{ReLU.} For $x\to+\infty$ only units with $w_h>0$ are active, giving the straight line
  \[
    \hat y\approx\Big(\sum_{w_h>0}v_hw_h\Big)x+\Big(\sum_{w_h>0}v_hb_h\Big)+c .
  \]
  A similar line holds for $x\to-\infty$ using the units with $w_h<0$.
\end{itemize}
So none of these networks can continue an oscillation. The cubic, by contrast, grows like $x^3$, which explains its very large error ($32$) on $[-2\pi,2\pi]$.

\subsection*{S9. Output layers for classification: softmax and cross-entropy}
For regression, as above, the output is linear and the loss is the mean squared error. For classification into $K$ classes, the lecture notes use the softmax output and cross-entropy loss.

Given output scores $z_1,\dots,z_K$ (logits), the \textbf{softmax} turns them into probabilities,
\[
  q_k=\frac{e^{z_k}}{\sum_je^{z_j}},\qquad q_k>0,\ \ \sum_kq_k=1 .
\]
With a target distribution $p$ (for example, one-hot for a known class), the \textbf{cross-entropy} is $L=-\sum_kp_k\ln q_k$. Using $\partial q_k/\partial z_j=q_k(\delta_{kj}-q_j)$,
\[
  \frac{\partial L}{\partial z_j}
  =-\sum_kp_k\frac{1}{q_k}\,q_k(\delta_{kj}-q_j)
  =-p_j+q_j\sum_kp_k=\boxed{\,q_j-p_j\,}.
\]
The gradient at the output is simply ``prediction minus target,'' the same form as the residual in regression.

\paragraph{Gibbs' inequality.} The KL divergence $\mathrm{KL}(p\|q)=\sum_kp_k\ln(p_k/q_k)$ is never negative. Using $\ln u\le u-1$ with $u=q_k/p_k$, summing over the classes with $p_k>0$:
\[
  -\mathrm{KL}(p\|q)=\sum_{p_k>0}p_k\ln\frac{q_k}{p_k}
  \le\sum_{p_k>0}(q_k-p_k)=\sum_{p_k>0}q_k-1\le0 .
\]
Hence $\mathrm{KL}\ge0$, with equality only when $q=p$. Because $\mathrm{KL}(p\|q)=H(p,q)-H(p)$ and $H(p)$ does not depend on the weights, minimizing cross-entropy and minimizing KL divergence are the same problem. (The lecture notes use $\log_2$; changing the base only rescales the loss by a constant.)

\subsection*{Check yourself}
\begin{enumerate}
  \item \textbf{Linear collapse.} For $W_1=[1,\,2]$, $\mathbf b_1=[0,\,1]$, $W_2=[3,\,-1]^{\T}$, $c=0.5$ and linear activations, compute the slope and intercept of the network by hand. Then check with the notebook starter cell.
  \item \textbf{Derivatives.} Compute $\sigma'(0)$, $\tanh'(0)$ and $\tanh'(3)$. Which of the two smooth activations passes more gradient through a unit that is far from $z=0$?
  \item \textbf{Backprop by hand.} One data point $x=1$, $y=1$, one tanh unit with $w=0.5$, $b=0$, $v=2$, $c=0$ and $N=1$. Compute $L$ and the four gradients $\partial L/\partial v$, $\partial L/\partial c$, $\partial L/\partial w$, $\partial L/\partial b$ by hand using S3, then verify them with the notebook.
  \item \textbf{Dead units.} For the ReLU networks of seeds 0 to 4, count the kinks inside $[-\pi,\pi]$. Does the count predict which seed has the lowest error? What other factors might matter?
  \item \textbf{Symmetry.} Why must the identical-unit network of S7 stay identical, even with Adam instead of plain gradient descent? (Hint: write the Adam update for unit $h$ and check what it depends on.)
\end{enumerate}

\end{document}
